\section{System Model and Threat Assumptions}
\label{sec:03}
\label{sec:system-model}

\subsection{System Model and Participating Entities}
\label{subsec:system-entities}

The present system retains identity-aware fingerprint and facial authorization from
the previously published VoteChain architecture \citep{rifat2026biometric}, but
separates that authorization from the candidate choice and the accepted private
ballot. The voter-facing client constructs an encrypted selection, while an
election-specific smart contract verifies its validity and accumulates ciphertexts
rather than plaintext candidate totals. The retained Semaphore-based registration
and election anchor is distinct from this private-voting eligibility path:
\texttt{VoteChainDirect} verifies membership against a frozen root of a separate
custom Poseidon commitment tree. The two structures are not interchangeable.
Biometric recognition performance belongs to the earlier study, not to the present
cryptographic evaluation.

The participating actors are the voter-facing web client and its privacy worker,
biometric authentication service, credential authorization service, election
administrator, blockchain contract and observers, five tally trustees, and an
external transcript auditor. Authentication identifies the voter; the credential
service manages registered secrets and authorizes protected credential release. The
client prepares the ballot and proof, whereas the contract checks the public
validity statement and records accepted encrypted ballots. The administrator
configures election state and submits recovered results. Trustees contribute to
threshold tally recovery, and the auditor can examine the corresponding evidence.
These responsibilities distinguish access to identities, choices, ciphertexts, and
decryption material; software separation alone does not demonstrate independent
institutional custody. Public observers can inspect transaction metadata,
proofs, nullifiers, ciphertext slots, and encrypted aggregates; election
administration alone does not confer access to credential-decryption keys or
trustee shares.

\subsection{Threat Model and Security Assumptions}
\label{subsec:system-security-objectives}

The threat model includes public observers of transactions, contract state,
senders, and nullifiers; submitters attempting ineligible, malformed, or repeated
ballots; and potentially compromised clients, credential custodians, or colluding
trustees. The intended acceptance guarantees are relative to the configured circuit
and contract checks. The Groth16 proof binds membership in the custom
eligible-credential tree, the election context, a one-hot candidate selection
within the active manifest, and correctness of the submitted encrypted slots
\citep{groth2016snark}. The contract rejects a reused nullifier and checks the
proof before accumulation. These mechanisms do not independently establish one
credential per civil identity, a fresh biometric match at submission, or integrity
of the client software.

Let $s$ denote a registered private voting-credential secret and $e$ the election
identifier. The implemented circuit derives the election-scoped nullifier as
follows:
\begin{equation}
  n_e = \mathsf{Poseidon}(d_{\mathrm{null}},s,e).
  \label{eq:system-nullifier}
\end{equation}
The nullifier is public, but the membership witness and selected candidate remain
private inputs to the proof. Here $d_{\mathrm{null}}$ is the configured domain
separator; Poseidon is a hash construction for proof circuits
\citep{grassi2021poseidon}. An observer cannot infer the plaintext candidate from
the nullifier alone. However, the credential service can recover a stored $s$
during authorization and recompute its $n_e$. An entity with this capability can
compare the result with recorded transactions and thereby link an identified voter
to an accepted ballot. Candidate-blind authorization therefore does not establish
universal identity unlinkability. A compromised voting client may also observe the
plaintext choice and credential before encryption.

Ballot confidentiality additionally depends on the election encryption key and
protected trustee shares. Three colluding trustees with sufficient shares could
decrypt individual published ciphertexts, and a three-of-five threshold does not
prove independent trustee administration. Groth16 ballot proofs are checked by the
contract, whereas distributed-key-generation and trustee partial-decryption
evidence are checked by an external verifier. The finalization transaction enforces
contract-state and arithmetic conditions and records totals with a transcript hash;
it does not verify the underlying decryption proofs. Independent result checking
requires the full transcript and the corresponding on-chain encrypted aggregates.
The present threat boundaries do not establish coercion resistance,
receipt-freeness, or production-election security.
